\section{Related Work}
\label{sec:related-work}

\subsection{XPIR and RLWE-based PIR}
XPIR~\cite{melchor2016xpir} is the closest backend reference. Its
homomorphic query and database evaluation support efficient polynomial
arithmetic, and its broader design includes aggregation and recursion.
Our scope retains the one-dimensional linear path, with a direct SKE
selector vector and explicit coefficient layout. This restricted path
does not exhaust XPIR's optimization space. The empirical conclusions below
concern concurrent same-profile repeated XPIR. No measured superiority
across all XPIR configurations is claimed. The pinned implementation~\cite{xpir_github} establishes
artifact provenance, not a substitute citation for the XPIR paper or
evidence that its native sampler realizes our formal distribution.

\subsection{Multi-value and Weighted Retrieval}
Hsu et al.~\cite[Secs.~III--V]{mulityquery} are the closest conceptual
prior work. Their group-homomorphic construction already supports an
encrypted weight at each selected position, encryptions of zero elsewhere,
a homomorphic weighted sum, and client-side recovery of two or more
values. They also discuss retrieval-count limits imposed by the plaintext
modulus and a privacy argument from the underlying encryption security.
Neither the multi-value goal nor that generic weighted architecture is
new here.

The positional weights \(K=(1,B,\ldots,B^{\alpha-1})\) are public and
target-independent. This instantiation replaces growing-weight
remainder recovery with ordinary radix digit extraction. Radix itself is
standard. The research increment is the precise connection between this
choice, record-to-polynomial encoding, exact capacity and the integer
no-wrap condition of the XPIR-style RLWE interface. The aggregate error
depends on all encoded database polynomials, whereas direct encryption of
a weight does not scale its encryption error by that weight. The formal
privacy reduction and native mapping concern this instantiation; generic
homomorphic privacy arguments and modulus limitations already occur in
the prior work. The conditional theorem alone establishes no concrete
security-admissible region. The separately tested native API frontier in
Section~\ref{sec:native-validation} concerns implementation feasibility and
finite functional execution, not a native noise or security certificate.

\subsection{Batch, SIMD, and Preprocessed PIR}
Angel et al.~\cite{angel2018pir} introduce SealPIR's compressed queries
and, separately, probabilistic batch codes (PBCs) for amortizing a batch.
The latter distribute encoded database items across buckets and use a
retrieval schedule, including cuckoo placement. SealPIR and the PBC
transformation are separate contributions in that same paper.
Vectorized BatchPIR~\cite{vectorized2023} uses vectorized RLWE homomorphic
encryption, bucketed batch retrieval, SIMD operations and ciphertext merging.
PIRANA~\cite{pirana2024} combines constant-weight index encoding with
homomorphic slot operations and a batch-code route; rotations and response
organization depend on the payload and configuration. Their setup,
evaluation keys and placement failure models matter to a comparison.

SimplePIR and DoublePIR are two protocols in the same
paper~\cite{henzinger2023one}. They use LWE-based linear processing with
database-dependent hints downloaded by the client; DoublePIR trades
additional processing for a smaller hint. Online work cannot be compared
to an entire fresh task while omitting this state and its lifecycle.
These protocols provide preprocessing context rather than the direct
weighted-packing precedent.

Respire~\cite[Secs.~2--4]{respire2024} directly studies small records and
batch queries. It uses subrings, ring switching, query compression and
response packing, and composes with PBCs for batch retrieval. Its model
includes reusable client query-key material and server database
preprocessing. SPIRIT~\cite[Secs.~2--5]{spirit2026} combines PBC scheduling
with LWE-to-RLWE conversion, shared key-expansion state and response
compression. Its hintless online interface does not mean absence of
server preprocessing or state. Thus neither short-record support,
small-batch access, lack of an offline client hint nor response packing
alone distinguishes the present work.

\begin{table}[t]
\centering
\caption{Mechanism summary, without a performance ranking. State and
assumptions differ across rows. The artifact provides an expanded mechanism comparison.}
\label{tab:mechanisms}
\small
\setlength{\tabcolsep}{3pt}
\begin{tabular}{@{}>{\raggedright\arraybackslash}p{2.0cm}>{\raggedright\arraybackslash}p{4.1cm}>{\raggedright\arraybackslash}p{4.7cm}@{}}
\toprule
Scheme & Retrieval mechanism & Relation to this study \\
\midrule
XPIR & Encrypted selection; linear evaluation; optional aggregation/recursion & Backend and empirical comparator; current path is flat. \\
Hsu et al. & Encrypted weights, weighted sum and integer recovery & Closest generic multi-value precedent; modulus/count limits already studied. \\
SealPIR / PBC & Compressed queries; separate bucketed batch transformation & Compression and placement are distinct mechanisms. \\
Vectorized BatchPIR & Bucket placement, SIMD processing and ciphertext merging & Batch alternative; not a local timing baseline. \\
PIRANA & Constant-weight selection, slots/rotations and batch codes & Alternative batch organization and payload tradeoffs. \\
SimplePIR / DoublePIR & LWE linear processing with client-downloaded hints & Preprocessing and amortization context. \\
Respire & Subring/ring-switching compression and response packing; batch codes & Direct short-record and small-batch precedent. \\
SPIRIT & PBCs, stateful ciphertext conversion and packed responses & Direct hintless-batch and response-packing precedent. \\
This construction & Several scalar digits per coefficient; unchanged linear ReplyGen & Restricted instantiation; capacity/noise/backend feasibility remains central. \\
\bottomrule
\end{tabular}
\end{table}


Radix packing here separates records by integer digit positions within
each coefficient, without bucket placement or SIMD rotations to recover
the targets. These architecture properties do not imply lower latency,
memory or failure probability. The contribution is the explicit
encoding/capacity/noise connection and its scoped evaluation within
XPIR, rather than query compression, offline hints or large-batch
optimization. External schemes provide mechanism and state context;
their published timings are not mixed with the local measurements.
